% TOP_PROBLEM: {"id": 30006940, "title": "Shub entropy conjecture", "category_id": 11, "category": {"id": 11, "name": "dynamical_systems", "display_name": "Dynamical Systems", "description": "Problems about long-term behavior of deterministic systems, Hamiltonian dynamics, and periodic orbits.", "slug": "dynamical-systems", "order_index": 11, "created_at": "2026-07-31T15:26:25.671Z"}, "status": "open"}
% FIRST_POSTED: 2026-09-27
% !TeX program = lualatex
% ATTEMPT_STATUS: unresolved
\documentclass[11pt]{article}
\usepackage[margin=25mm]{geometry}
\usepackage{fontspec}
\setmainfont{Latin Modern Roman}
\usepackage{amsmath,amssymb,mathtools}
\usepackage{unicode-math}
\setmathfont{Latin Modern Math}
\usepackage{xurl,hyperref}
\hypersetup{colorlinks=true,urlcolor=blue}
\setcounter{secnumdepth}{0}
\setlength{\parindent}{0pt}
\setlength{\parskip}{5pt}
\setlength{\emergencystretch}{3em}
\begin{document}
\section*{\texorpdfstring{Shub entropy conjecture}{Shub entropy conjecture}}\label{30006940}
\subsection{Definitions and mathematical statement}
Let $f:M\to M$ be a $C^1$ self-map of a compact smooth manifold. For a compatible metric set $d_N(x,y)=\max_{0\le j<N}d(f^jx,f^jy)$ and let $s_N(\delta)$ be the maximum cardinality of a $\delta$-separated set for $d_N$. The topological entropy is
\[
 h_{\rm top}(f)=\lim_{\delta\downarrow0}\limsup_{N\to\infty}
 \frac1N\log s_N(\delta).
\]
Let $f_*$ act on the finite-dimensional graded real homology $H_*(M;{\mathbb{R}})$, and let $\rho(f_*)$ be its spectral radius, the largest absolute eigenvalue. Shub's assertion is
\[
 h_{\rm top}(f)\ge\log\rho(f_*).
\]
Invertibility is not assumed. The differentiability class $C^1$ is part of the target; replacing it with a stronger smoothness hypothesis narrows the question.
\par\smallskip\textit{Formulation and concept references:} \href{https://arxiv.org/pdf/2112.14181}{[S1]}.

\subsection{Short English statement}
Must the dynamical entropy of a continuously differentiable map dominate the growth rate it induces on homology?
\subsection{Research attempt}
\paragraph{Process, scope, and first missing step.}
This attempt was generated by GPT-6 Astra Ultra on 27 September 2026. We compute the conjectured lower bound exactly for diagonal affine maps of tori, then examine two routes toward a general $C^1$ map: growth of representative cycles and approximation by smooth maps. The computations are restricted results, not a new proof of Shub's conjecture. The missing implication is a sufficiently sharp conversion from homological growth to the number of distinguishable orbits at a fixed spatial resolution.

The introduction of [S1] records the smooth case and the first-homology lower bound, while retaining the general $C^1$ problem. Its extension involving first Cech cohomology does not replace the full graded homology in this question. A theorem for diffeomorphisms away from homoclinic tangencies [E1] has additional dynamical hypotheses and invertibility. We do not impose them on the catalogue statement. The arguments below use elementary calculations independently of these stronger theorems.

\paragraph{A family with an exact answer in every dimension.}
Let $M=\mathbb T^d=\mathbb R^d/\mathbb Z^d$ and
\[
 f(x_1,\ldots,x_d)=(q_1x_1+c_1,\ldots,q_dx_d+c_d)
       \pmod {\mathbb Z^d},\qquad q_i\in\mathbb Z.
 \tag{445.1}
\]
Use the maximum of the circle distances. Negative degrees, zero degrees, and arbitrary translations are allowed. Write $Q_i=\max(1,|q_i|)$. We prove
\[
 h_{\rm top}(f)=\sum_{i=1}^d\log Q_i
              =\log\rho(f_*).
 \tag{445.2}
\]
For the upper bound, partition the $i$th initial circle into at most $C_\delta Q_i^{N-1}$ half-open intervals of length less than $\delta Q_i^{-(N-1)}$. If two initial points belong to the same product box, their coordinate differences after $j<N$ iterates are at most $|q_i|^j$ times their initial differences, before taking circle distance. This is less than $\delta$ for every coordinate. For $q_i=0$, the difference vanishes after the first iterate; for $|q_i|=1$, it stays unchanged. Each box therefore contains at most one member of a $\delta$-separated set. Consequently
\[
 s_N(\delta)\le C_\delta^d\prod_i Q_i^{N-1}.
 \tag{445.3}
\]
The constants depend on the fixed scale and dimension, not on $N$.

For the lower bound, in each coordinate with $Q_i\ge2$ use the grid $\{a/Q_i^{N-1}:0\le a<Q_i^{N-1}\}$, and use a single initial value in every other coordinate. Two distinct grid entries differ by $k/Q_i^{N-1}$ with $0<|k|<Q_i^{N-1}$. Write $k=Q_i^r\ell$, with $Q_i\nmid\ell$ and $r\le N-2$. After $N-2-r$ iterates their difference modulo one is, up to sign, $\ell/Q_i$. Its circle distance from zero is at least $1/Q_i$. Translations cancel when differences of iterates are taken.

Thus the product grid, of size $\prod_i Q_i^{N-1}$, is separated for any fixed positive $\delta<\min_{Q_i\ge2}1/Q_i$. If there are no expanding coordinates, entropy is nonnegative and (445.3) already gives zero. Taking the limits in the definition proves the first equality in (445.2). This argument counts distinguishable orbits directly; it does not infer entropy from a derivative norm alone.

\paragraph{Checking the graded homology, including degree zero.}
The real homology of a torus is the exterior algebra on $H_1(\mathbb T^d;\mathbb R)\cong\mathbb R^d$. Translation is homotopic to the identity. On the wedge vector indexed by a subset $I\subseteq\{1,\ldots,d\}$, the induced map has eigenvalue $\prod_{i\in I}q_i$. The empty subset contributes the eigenvalue $1$ on $H_0$. Therefore
\[
 \rho(f_*)=\max_I\left|\prod_{i\in I}q_i\right|
          =\prod_i\max(1,|q_i|).
 \tag{445.4}
\]
Zero eigenvalues cause no logarithmic ambiguity because the $H_0$ eigenvalue remains. Formula (445.4) also shows why merely quoting a first-homology theorem can lose information. For $q_1=\cdots=q_d=2$, the first-homology radius is $2$, whereas the graded radius is $2^d$. The required lower bound is $d\log2$, not $\log2$. Applying the same first-homology result to iterates only multiplies both of those logarithms by the iterate number; it does not supply the missing factor $d$.

\paragraph{What an eigenclass really proves about volume.}
Suppose a positive-degree cohomology class $[\omega]\in H^k(M;\mathbb C)$ satisfies $f^*[\omega]=\lambda[\omega]$, with $\lambda\ne0$. Choose a smooth closed differential-form representative and a finite piecewise smooth real $k$-cycle $c$ with $\int_c\omega\ne0$. Complex coefficients can be handled by the real and imaginary parts. Cohomological naturality gives
\[
 \int_c(f^n)^*\omega=\lambda^n\int_c\omega.
 \tag{445.5}
\]
With respect to a fixed Riemannian metric, the comass of $\omega$ is bounded on compact $M$. If $V_n(c)$ denotes the sum of the parametrized $k$-volumes of the images of the simplices of $c$, counted with absolute coefficients and multiplicity, then
\[
 V_n(c)\ge
 \frac{|\int_c\omega|}{\|\omega\|_{\rm comass}}|\lambda|^n.
 \tag{445.6}
\]
This is a genuine geometric consequence of the eigenvalue. It neither bounds the volume of the image as a set nor counts separated orbits. Self-overlap can make parametrized volume substantially larger than the measure of the image.

At a fixed resolution, large parametrized volume alone gives no lower bound on a covering number. For instance a curve can traverse a small circle arbitrarily many times while its image lies inside a single prescribed ball. This example is not claimed to arise from iterates of a fixed map with the required homology. Its purpose is to expose the invalid general inference from the left side of (445.6) to separated-point growth. A successful iteration argument must use dynamical information to control multiplicity and small-scale folding. The displayed volume inequality by itself contains no such control.

\paragraph{A smooth-approximation route and its exact continuity requirement.}
Let smooth maps $f_j$ converge to $f$ in $C^1$. For sufficiently close maps the homotopy class, hence the induced homology map, is unchanged. On a compact manifold this follows by joining nearby image points inside a common convex normal neighborhood. The smooth entropy theorem cited in [S1] would give
\[
 h_{\rm top}(f_j)\ge\log\rho(f_*).
 \tag{445.7}
\]
To pass this lower bound to $f$, the useful inequality would be
\[
 h_{\rm top}(f)\ge\limsup_{j\to\infty}h_{\rm top}(f_j).
 \tag{445.8}
\]
This is upper semicontinuity at $f$. Lower semicontinuity has the opposite direction and would not complete the proof. Nor does convergence of a fixed finite collection of iterates imply (445.8): the entropy definition first permits orbit lengths tending to infinity, and the required closeness of $f_j^N$ to $f^N$ need not be uniform in $N$.

The special dynamical class in [E1] is an example where additional control is available. The present attempt proves no corresponding statement for every $C^1$ map. Uniform bounds on a finite number of derivatives of approximants also cannot be asserted merely from $C^1$ convergence. This is the precise point where replacing the hypothesis by $C^\infty$ silently changes the problem.

\paragraph{Verification and outcome.}
Exact finite checks enumerate the circle grids used above for small positive and negative degrees, including zero and unit degrees, and verify the exterior-product maximum for integer diagonal matrices. The proofs of (445.2)--(445.6) are written separately because finite samples cannot verify entropy limits or arbitrary smooth cycles. Source hypotheses were checked against the stated differentiability class and the absence of an invertibility assumption.

\textbf{UNRESOLVED.} The conjecture holds with equality for the diagonal affine family (445.1). The general volume argument stops before a fixed-scale orbit-counting estimate, and the approximation argument stops at (445.8). Neither missing statement is established for arbitrary $C^1$ self-maps.

\subsection{Sources}
\begin{itemize}
\item[S1]Hernandez-Corbato, Nieves-Rivera, Ruiz del Portal, and Sanchez-Gabites, Integration in Cech theories and a bound on entropy, Introduction.
\url{https://arxiv.org/pdf/2112.14181}.
\textit{Formulation source.}
\item[E1]G. Liao, M. Viana, and J. Yang, \emph{The Entropy Conjecture for Diffeomorphisms away from Tangencies}, arXiv:1012.0514.
\url{https://arxiv.org/abs/1012.0514}.
\textit{Restricted differentiable-dynamics theorem; consulted 27 September 2026.}
\end{itemize}
\end{document}
